\documentclass[12pt]{article}
\usepackage[utf8]{inputenc}
\usepackage[T1]{fontenc}
\usepackage{hyperref}
\title{Panim Language Resources for K-12 Curriculum Development}
\author{Compiled by Claude Sonnet 5.5 (claude-sonnet-5-5) for Andr\'as Kornai}
\date{}
\begin{document}
\maketitle
\section{Introduction}
Panim (ISO 639-3 pnr) is a Madang (Trans-New Guinea) language of Madang Province, Papua New Guinea, spoken in the village of Panim. Children now learn Tok Pisin. Speaker estimates conflict: Wikipedia gives 12 (2019), another source about 10 elderly speakers plus heritage speakers, and Living Tongues roughly 400 people in the village.
\section{Primary Educational Resources}
No school materials were found. The talking dictionary below could support classroom use.
\section{Linguistic and Reference Materials}
\begin{itemize}
\item Panim Talking Dictionary (Swarthmore, with Living Tongues Institute and Don Daniels): about page. \url{https://talkingdictionary.swarthmore.edu/panim/about}
\item Wikipedia, Panim language. \url{https://en.wikipedia.org/wiki/Panim_language}
\item PARADISEC catalogue, Daniels Panim notebooks, collection DD1. \url{https://catalog.paradisec.org.au/repository/DD1/001} and \url{https://catalog.paradisec.org.au/repository/DD1/002}
\end{itemize}
\section{Religious and Community Resources}
\begin{itemize}
\item Joshua Project, Panim in Papua New Guinea (profile). \url{https://joshuaproject.net/index.php/people_groups/14297}
\end{itemize}
\section{Digital and Audio Resources}
\begin{itemize}
\item Living Tongues Institute, Madang Province page, describing the Panim documentation work. \url{https://livingtongues.org/madang-province/}
\end{itemize}
\section{Curriculum Development Context}
Panim has a talking dictionary and an archived fieldwork collection, which is more than most languages in this batch have. Transmission to children has stopped, so a curriculum would be a revitalization effort built on the few elderly speakers and the existing documentation.
\end{document}